\section{How to read this manual}

This manual is compiled \emph{with} the toolkit it documents: its class is the
shipped \file{CourseDocument}, its packages are the shipped ones, and every
rendered example below is typeset by the real machinery at build time. If the
manual builds, the toolkit works.

Examples appear in boxes like this one -- the source on top is the \emph{same
text} that produced the rendered result underneath, so the two can never drift
apart:

\begin{manexample}
A limit shorthand from MathStuff:
$\inflim \frac{1}{n} = 0$.
\end{manexample}

Three things to keep in mind while reading:

\begin{itemize}
  \item \textbf{Every content layer is switched on.} This manual forces
    solutions, hints and (locally) instructor notes to display, because its job
    is to show them. In a real course document these layers are \emph{off}; the
    build script turns them on per view, producing separate student, solutions,
    hints and instructor PDFs from one source. Part 1 explains that model.
  \item \textbf{Counters are live.} A rendered \env{problem} steps the real
    problem counter, so example problems number consecutively through the
    manual, exactly as they would in a course document.
  \item \textbf{Solid mode is active.} Boxes here occupy their natural height.
    The default \emph{overlay} mode -- solution boxes of zero height that paint
    into reserved whitespace so every view has identical page breaks -- has its
    own chapter, and one live demonstration there.
\end{itemize}

Where to start: section~\ref{sec:quickstart} for the shape of a course
repo, then whichever workflow you are sitting down to do. The reference
(Part~\ref{part:reference}) and the architecture notes
(Part~\ref{part:architecture}) are for looking things up later.
